\section{Общие утверждения из прошлых курсов}

\begin{definition}[Функция распределения]
\leavevmode

$F_\xi(x) = P(\xi < x)$ --- функция распределения вероятности.
Для непрерывного случая:
$$ F_\xi(x) = \int_{-\infty}^{x} f(t) dt $$
\end{definition}

\begin{definition}[Математическое ожидание и дисперсия]
\leavevmode

Математическое ожидание (мат. ожидание):
\begin{gather*}
E\xi = \sum_{i=1}^{n} p_i x_i \\
E\xi = \int_{-\infty}^{+\infty} x f(x) \, dx \\
E\phi(\xi) = \int_{-\infty}^{+\infty} \phi(x) f(x) \, dx
\end{gather*}
Дисперсия:
$$ D\xi = E\left((\xi - E\xi)^2\right) = E(\xi^2) - 2(E\xi)^2 + E(\xi)^2 = E(\xi^2) - (E(\xi))^2 $$
\end{definition}

\begin{definition}[Нормальное распределение]
\leavevmode

Плотность нормального распределения:
\begin{gather*}
\xi \sim N(\mu, \sigma)\\
f(x) = \frac{1}{\sqrt{2\pi\sigma^2}} \cdot e^{-\frac{(x-\mu)^2}{2\sigma^2}}, \quad -\infty < x < \infty 
\end{gather*}

При $\mu=0, \sigma=1$ ~--- стандартное нормальное распределение:
\begin{gather*}
\xi \sim N(0, 1) \\
f(x) = \frac{1}{\sqrt{2\pi}} \cdot e^{-\frac{x^2}{2}}
\end{gather*}


\end{definition}
